%----------------------------------------------------------------------------------------
% TRANSFORMATIES VAN GRAFIEKEN. Alle opgaven zijn eigen werk.
%
% Opbouw: Start, leerstof (theorie, werkwijze of uitleg, voorbeeld, korte oefening), en op een
% nieuwe pagina de Opdrachten (door elkaar, stijgende moeilijkheid), Koppel en Test jezelf.
% Labels van dit hoofdstuk beginnen met h2-.
%----------------------------------------------------------------------------------------

\mychapter{2}{Transformaties van grafieken}

\lkr{Leerplan B+S'': LPD 9-11 (algemene functies), het deel over transformaties van grafieken. In de cursussen over exponentiële, goniometrische en rationale functies komen dezelfde transformaties terug: daar volstaat dan een verwijzing naar dit hoofdstuk. Laat leerlingen telkens één bijzonder punt volgen (een top, een beginpunt): dat maakt de transformatie zichtbaar en controleerbaar.}

Veel functies zijn een bekende functie in een andere vorm: verschoven, uitgerekt of gespiegeld. Wie de grafieken van een handvol basisfuncties kent, kan zo heel wat grafieken schetsen zonder tabel. Het idee is telkens hetzelfde: een grafiek is een verzameling punten (hoofdstuk 1), en een kleine verandering in het voorschrift verplaatst elk punt op dezelfde manier. We zoeken dus uit waar elk punt $(x, y)$ terechtkomt.

\begin{oefeningen}[Start]
\oef{h2-start-1}

\oef{h2-verken}
\end{oefeningen}

%========================================================================================
\section{Basisfuncties}

Deze grafieken ken je best uit het hoofd, met hun domein en bereik.

\figuren{%
\begin{tikzpicture}[x=.36cm,y=.36cm,font=\scriptsize]
  \roosterleeg{-3}{1}{3}{-3}{1}{3}
  \kromme{-1.73}{1.73}{(\x)^2}
  \node[below,text=black] at (0,-3.2) {$x^2$};
\end{tikzpicture}\hspace{3mm}
\begin{tikzpicture}[x=.36cm,y=.36cm,font=\scriptsize]
  \roosterleeg{-3}{1}{3}{-3}{1}{3}
  \kromme{-1.44}{1.44}{(\x)^3}
  \node[below,text=black] at (0,-3.2) {$x^3$};
\end{tikzpicture}\hspace{3mm}
\begin{tikzpicture}[x=.36cm,y=.36cm,font=\scriptsize]
  \roosterleeg{-3}{1}{3}{-3}{1}{3}
  \kromme{0}{3}{sqrt(\x)}
  \vol{(0,0)}
  \node[below,text=black] at (0,-3.2) {$\sqrt{x}$};
\end{tikzpicture}\hspace{3mm}
\begin{tikzpicture}[x=.36cm,y=.36cm,font=\scriptsize]
  \roosterleeg{-3}{1}{3}{-3}{1}{3}
  \draw[accent,thick] (-3,3) -- (0,0) -- (3,3);
  \node[below,text=black] at (0,-3.2) {$|x|$};
\end{tikzpicture}\hspace{3mm}
\begin{tikzpicture}[x=.36cm,y=.36cm,font=\scriptsize]
  \roosterleeg{-3}{1}{3}{-3}{1}{3}
  \kromme{.34}{3}{1/\x} \kromme{-3}{-.34}{1/\x}
  \node[below,text=black] at (0,-3.2) {$\frac1x$};
\end{tikzpicture}}

\begin{theorie}
\renewcommand{\arraystretch}{1.25}
\begin{tabular}{@{}l@{\hspace{8mm}}ccccc@{}}
$f(x)$ & $\formule{x^2}$ & $\formule{x^3}$ & $\formule{\sqrt{x}}$ & $\formule{|x|}$ & $\formule{\frac1x}$ \\
$\dom f$ & $\R$ & $\R$ & $\left[0, +\infty\right[$ & $\R$ & $\R \setminus \{0\}$ \\
$\ber f$ & $\left[0, +\infty\right[$ & $\R$ & $\left[0, +\infty\right[$ & $\left[0, +\infty\right[$ & $\R \setminus \{0\}$ \\
symmetrie & even & oneven & geen & even & oneven
\end{tabular}
\end{theorie}

Je hoeft de tabel niet blind te leren: elke kolom volgt uit het voorschrift, met de begrippen uit hoofdstuk 1. Voor het bereik zeg je telkens welke waarden mogelijk zijn, en dat elke zulke waarde ook bereikt wordt.

\begin{itemize}
\item $x^2$: een kwadraat is nooit negatief, en elke $y \geqslant 0$ is het kwadraat van $\sqrt{y}$. Even, want $(-x)^2 = x^2$.
\item $x^3$: elke $y$ is de derdemacht van $\sqrt[3]{y}$, dus het bereik is $\R$. Oneven, want $(-x)^3 = -x^3$.
\item $\sqrt{x}$: bestaat alleen voor $x \geqslant 0$ en is nooit negatief; elke $y \geqslant 0$ is $\sqrt{y^2}$. Het domein is niet symmetrisch, dus geen symmetrie.
\item $|x|$: nooit negatief, en $|y| = y$ voor $y \geqslant 0$. Even, want $|-x| = |x|$.
\item $\frac1x$: bestaat niet in $0$, en is nooit $0$; elke $y \neq 0$ is het beeld van $\frac1y$. Oneven, want $\frac{1}{-x} = -\frac1x$.
\end{itemize}

\begin{oefeningen}
\oef{h2-basis}
\end{oefeningen}

%========================================================================================
\section{Verschuiven}

In de Verken was $g(x) = (x - 1)^2$ de parabool $y = x^2$, één eenheid naar rechts geschoven. We bekijken eerst een verschuiving in de hoogte, dan een opzij, en ten slotte beide samen.

\subsection{Verticaal verschuiven}

Tel je bij elke functiewaarde hetzelfde getal $d$ op, dan komt elk punt van de grafiek $d$ hoger te liggen.

\begin{theorie}
\begin{stelling}{Verticale verschuiving}
De grafiek van $\formule{g(x) = f(x) + d}$ is die van $f$, verschoven over $d$ naar boven (naar onder als $d < 0$): \quad $\formule{(x, y) \mapsto (x, y + d)}$
\end{stelling}

\begin{bewijs}
$g$ en $f$ hebben hetzelfde domein. Een punt $(x, y)$ ligt op de grafiek van $g$ als en slechts als
\bewijsregel{y = f(x) + d \iff y - d = f(x) ,}
dus als en slechts als $(x, y - d)$ op de grafiek van $f$ ligt. De punten van $g$ zijn dus precies de punten van $f$ met $d$ opgeteld bij de $y$-coördinaat.
\end{bewijs}
\end{theorie}

Het domein verandert niet, het bereik schuift $d$ mee op.

\voorbeeld{}{Schets de grafiek van $g(x) = |x| - 2$, en geef het bereik en de nulwaarden.\label{vb:h2-vertverschuiven}}{%
\naast{%
\uitleg{Dit is $|x|$ met $d = -2$. Volg de knik $(0, 0)$.}
\begin{stappenvrij}{4.2cm}{4.6cm}
(0, 0) \mapsto (0, -2) & De knik zakt 2 eenheden. \\
\ber g = \left[-2, +\infty\right[ & Het bereik $\left[0, +\infty\right[$ zakt mee. \\
|x| - 2 = 0 \iff x = \pm 2 & De nulwaarden.
\end{stappenvrij}}{%
\begin{tikzpicture}[x=.45cm,y=.45cm,font=\scriptsize]
  \rooster{-3}{1}{3}{-2}{1}{2}
  \draw[grijs,dashed,thick] (-2.4,2.4) -- (0,0) -- (2.4,2.4);
  \draw[accent,thick] (-3,1) -- (0,-2) -- (3,1);
  \vol{(0,-2)}
\end{tikzpicture}}}

\begin{oefeningen}
\oef{h2-vert-kort}
\end{oefeningen}

\subsection{Horizontaal verschuiven}

Bij $g(x) = f(x - c)$ verandert iets aan de $x$. Wat $f$ doet in $0$, doet $g$ in $x = c$, want $g(c) = f(0)$. Alles gebeurt dus $c$ eenheden later.

\begin{theorie}
\begin{stelling}{Horizontale verschuiving}
De grafiek van $\formule{g(x) = f(x - c)}$ is die van $f$, verschoven over $c$ naar rechts (naar links als $c < 0$): \quad $\formule{(x, y) \mapsto (x + c, y)}$
\end{stelling}

\begin{bewijs}
Een punt $(x, y)$ ligt op de grafiek van $g$ als en slechts als
\bewijsregel{y = f(x - c) ,}
dus als en slechts als $(x - c, y)$ op de grafiek van $f$ ligt. Noem $u = x - c$. Bij elk punt $(u, y)$ van $f$ hoort dus het punt $(u + c, y)$ van $g$, en omgekeerd.
\end{bewijs}
\end{theorie}

Let op het teken: in $f(x - 3)$ schuift de grafiek naar \textit{rechts}, in $f(x + 3) = f\bigl(x - (-3)\bigr)$ naar links. Nu verandert het domein mee, en blijft het bereik hetzelfde.

\voorbeeld{}{Schets de grafiek van $g(x) = \dfrac{1}{x - 1}$, en geef het domein.\label{vb:h2-horverschuiven}}{%
\naast{%
\uitleg{Dit is $\frac1x$ met $c = 1$. Volg twee punten, en de verticale rechte waar de functie niet bestaat.}
\begin{stappenvrij}{4.2cm}{4.6cm}
(1, 1) \mapsto (2, 1) & 1 naar rechts. \\
(-1, -1) \mapsto (0, -1) & Controle: $g(0) = -1$. \\
x = 0 \ \text{wordt}\ x = 1 & Daar bestaat $g$ niet. \\
\dom g = \R \setminus \{1\} & Het bereik blijft $\R \setminus \{0\}$.
\end{stappenvrij}}{%
\begin{tikzpicture}[x=.45cm,y=.45cm,font=\scriptsize]
  \rooster{-2}{1}{4}{-3}{1}{3}
  \draw[grijs,dashed] (1,-3) -- (1,3.3);
  \kromme{1.34}{4}{1/(\x-1)} \kromme{-2}{.66}{1/(\x-1)}
  \vol{(2,1)} \vol{(0,-1)}
\end{tikzpicture}}}

\begin{oefeningen}
\oef{h2-hor-kort}
\end{oefeningen}

\subsection{Beide samen}

Bij $g(x) = f(x - c) + d$ pas je beide stellingen na elkaar toe: eerst $c$ opzij, dan $d$ omhoog. Een punt $(x, y)$ gaat zo naar $(x + c, y + d)$. De volgorde maakt hier niet uit, want de ene verschuiving verandert alleen de $x$, de andere alleen de $y$.

\voorbeeld{}{Schets de grafiek van $g(x) = \sqrt{x + 3} - 1$, en geef het domein en het bereik.\label{vb:h2-verschuiven}}{%
\naast{%
\uitleg{Schrijf $x + 3$ als $x - (-3)$: dat is $c = -3$, en $d = -1$. Volg het beginpunt $(0, 0)$ van $\sqrt{x}$.}
\begin{stappenvrij}{4.3cm}{3.3cm}
c = -3,\ d = -1 & 3 naar links, 1 naar onder. \\
(0, 0) \mapsto (-3, -1) & Het beginpunt. \\
\dom g = \left[-3, +\infty\right[ & Ook 3 naar links. \\
\ber g = \left[-1, +\infty\right[ & Ook 1 naar onder.
\end{stappenvrij}}{%
\begin{tikzpicture}[x=.5cm,y=.5cm,font=\scriptsize]
  \rooster{-3}{1}{4}{-1}{1}{3}
  \draw[grijs,dashed,thick] plot[domain=0:4,samples=60,smooth] (\x,{sqrt(\x)});
  \kromme{-3}{4}{sqrt(\x+3)-1}
  \vol{(-3,-1)}
\end{tikzpicture}}}

\nujij{Waar ligt de top van de parabool $y = (x - 2)^2 + 1$?}{De top $(0, 0)$ van $y = x^2$ schuift 2 naar rechts en 1 naar boven: $(2, 1)$.}

\begin{oefeningen}
\oef{h2-verschuiven}

\needspace{8\baselineskip}
\oef{h2-verschuiven-grafiek}
\end{oefeningen}

%========================================================================================
\section{Uitrekken en spiegelen}

Een verschuiving telt iets op. Vermenigvuldig je met een getal, dan wordt de grafiek uitgerekt of samengedrukt. Ook hier hangt het ervan af of je dat buiten $f$ doet of binnen.

\subsection{Verticaal uitrekken}

\begin{theorie}
\begin{stelling}{Verticaal uitrekken}
Voor $a > 0$ is de grafiek van $\formule{g(x) = a \cdot f(x)}$ die van $f$, verticaal uitgerekt met factor $a$: \quad $\formule{(x, y) \mapsto (x, ay)}$
\end{stelling}

\begin{bewijs}
$g$ en $f$ hebben hetzelfde domein. Een punt $(x, y)$ ligt op de grafiek van $g$ als en slechts als
\bewijsregel{y = a f(x) \iff \frac{y}{a} = f(x) ,}
want $a \neq 0$. Dat is: als en slechts als $\bigl(x, \frac{y}{a}\bigr)$ op de grafiek van $f$ ligt. Bij elk punt $(x, v)$ van $f$ hoort dus het punt $(x, av)$ van $g$, en omgekeerd.
\end{bewijs}
\end{theorie}

Is $a > 1$, dan wordt de grafiek steiler; is $0 < a < 1$, dan wordt ze platter. Punten op de $x$-as blijven liggen, want $a \cdot 0 = 0$: de nulwaarden veranderen niet.

\voorbeeld{}{Schets de grafiek van $g(x) = \frac12 x^2$.\label{vb:h2-vertrekken}}{%
\naast{%
\uitleg{Dit is $x^2$ met $a = \frac12$: elke hoogte halveert. Volg enkele punten van $x^2$.}
\begin{stappenvrij}{4.2cm}{4.4cm}
(0, 0) \mapsto (0, 0) & De top blijft liggen. \\
(\pm 2, 4) \mapsto (\pm 2, 2) & Half zo hoog. \\
(\pm 3, 9) \mapsto (\pm 3, \tfrac92) & Controle: $\frac12 \cdot 9 = \frac92$.
\end{stappenvrij}
De parabool is breder geworden, maar ze is niet opzij geschoven.}{%
\begin{tikzpicture}[x=.45cm,y=.4cm,font=\scriptsize]
  \rooster{-3}{1}{3}{0}{1}{5}
  \draw[grijs,dashed,thick] plot[domain=-2.2:2.2,samples=60,smooth] (\x,{(\x)^2});
  \kromme{-3.1}{3.1}{(\x)^2/2}
  \vol{(-2,2)} \vol{(2,2)}
\end{tikzpicture}}}

\begin{oefeningen}
\oef{h2-vertrekken-kort}
\end{oefeningen}

\subsection{Horizontaal uitrekken}

\begin{theorie}
\begin{stelling}{Horizontaal uitrekken}
Voor $b > 0$ is de grafiek van $\formule{h(x) = f(bx)}$ die van $f$, horizontaal uitgerekt met factor $\frac1b$: \quad $\formule{(x, y) \mapsto \left(\frac{x}{b}, y\right)}$
\end{stelling}

\begin{bewijs}
Een punt $(x, y)$ ligt op de grafiek van $h$ als en slechts als
\bewijsregel{y = f(bx) ,}
dus als en slechts als $(bx, y)$ op de grafiek van $f$ ligt. Noem $u = bx$, dus $x = \frac{u}{b}$. Bij elk punt $(u, y)$ van $f$ hoort het punt $\bigl(\frac{u}{b}, y\bigr)$ van $h$, en omgekeerd.
\end{bewijs}
\end{theorie}

Ook hier werkt een factor binnen de haakjes \textit{omgekeerd}: $f(2x)$ maakt de grafiek twee keer smaller, want $h(1) = f(2)$: wat $f$ in $2$ doet, doet $h$ al in $1$. Punten op de $y$-as blijven liggen, want $\frac0b = 0$.

\voorbeeld{}{De parabool $f(x) = x^2 - 4x$ heeft de nulwaarden $0$ en $4$ en de top $(2, -4)$. Bepaal de nulwaarden en de top van $h(x) = f(2x)$, en controleer met het voorschrift.\label{vb:h2-horrekken}}{%
\naast{%
\uitleg{Met $b = 2$ halveren alle $x$-coördinaten, de $y$-coördinaten blijven.}
\begin{stappenvrij}{4.2cm}{4.4cm}
0 \mapsto 0,\ 4 \mapsto 2 & Nulwaarden van $h$: $0$ en $2$. \\
(2, -4) \mapsto (1, -4) & De top van $h$. \\
h(x) = (2x)^2 - 4 \cdot 2x & Controle met het voorschrift: \\
\phantom{h(x)} = 4x(x - 2) & nulwaarden $0$ en $2$, \\
h(1) = 4 - 8 = -4 & en de top klopt.
\end{stappenvrij}}{%
\begin{tikzpicture}[x=.45cm,y=.4cm,font=\scriptsize]
  \rooster{-1}{1}{5}{-4}{1}{2}
  \draw[grijs,dashed,thick] plot[domain=-.45:4.45,samples=60,smooth] (\x,{(\x)^2-4*\x});
  \kromme{-.4}{2.4}{4*(\x)^2-8*\x}
  \vol{(1,-4)} \fill[grijs] (2,-4) circle (2.2pt);
\end{tikzpicture}}}

\begin{oefeningen}
\oef{h2-horrekken-kort}
\end{oefeningen}

\subsection{Verticaal of horizontaal?}

Het verschil tussen $a \cdot f(x)$ en $f(bx)$ zie je het best naast elkaar, met dezelfde grafiek en dezelfde factor.

\voorbeeld{}{De grafiek van $f$ is de gebroken lijn door $(-2, 0)$, $(0, 2)$, $(2, 0)$ en $(4, -2)$. Teken de grafieken van $g(x) = 2f(x)$ en $h(x) = f(2x)$.\label{vb:h2-rekken}}{%
\uitleg{Transformeer de vier hoekpunten en verbind ze. Bij $g$ verdubbelen de $y$-waarden, bij $h$ halveren de $x$-waarden.}
\begin{stappenbreed}
g\colon\ (-2, 0),\ (0, 4),\ (2, 0),\ (4, -4) & $(x, y) \mapsto (x, 2y)$: twee keer zo hoog. \\
h\colon\ (-1, 0),\ (0, 2),\ (1, 0),\ (2, -2) & $(x, y) \mapsto (\frac{x}{2}, y)$: twee keer zo smal.
\end{stappenbreed}
\figuren{%
\begin{tikzpicture}[x=.42cm,y=.42cm,font=\scriptsize]
  \rooster{-2}{1}{4}{-4}{1}{4}
  \draw[grijs,dashed,thick] (-2,0) -- (0,2) -- (2,0) -- (4,-2);
  \draw[accent,thick] (-2,0) -- (0,4) -- (2,0) -- (4,-4);
  \node[right,text=accent] at (.4,3.6) {$g$};
\end{tikzpicture}\hspace{1cm}
\begin{tikzpicture}[x=.42cm,y=.42cm,font=\scriptsize]
  \rooster{-2}{1}{4}{-4}{1}{4}
  \draw[grijs,dashed,thick] (-2,0) -- (0,2) -- (2,0) -- (4,-2);
  \draw[accent,thick] (-1,0) -- (0,2) -- (1,0) -- (2,-2);
  \node[right,text=accent] at (1.6,-2.3) {$h$};
\end{tikzpicture}}
Bij $g$ blijven de nulwaarden liggen, bij $h$ het snijpunt $(0, 2)$ met de $y$-as.}

\subsection{Spiegelen om de $x$-as}

Wat gebeurt er als de factor negatief is? Neem eerst $a = -1$: elke functiewaarde verandert van teken.

\begin{theorie}
\begin{stelling}{Spiegelen om de $x$-as}
De grafiek van $\formule{g(x) = -f(x)}$ is het spiegelbeeld van die van $f$ om de $x$-as: \quad $\formule{(x, y) \mapsto (x, -y)}$
\end{stelling}

\begin{bewijs}
Het bewijs van verticaal uitrekken gebruikt alleen dat $a \neq 0$. Met $a = -1$ geeft het: bij elk punt $(x, v)$ van $f$ hoort het punt $(x, -v)$ van $g$. Dat is het spiegelbeeld om de $x$-as.
\end{bewijs}
\end{theorie}

Een negatieve factor combineert dus twee transformaties: $-3f(x) = -\bigl(3f(x)\bigr)$ is verticaal uitrekken met factor $3$ en daarna spiegelen om de $x$-as.

\voorbeeld{}{Geef het domein en het bereik van $k(x) = -\sqrt{x}$.\label{vb:h2-spiegelx}}{%
\naast{%
\uitleg{Spiegel $\sqrt{x}$ om de $x$-as: de $x$-coördinaten blijven, de $y$-coördinaten veranderen van teken.}
\begin{stappenvrij}{4.2cm}{4.2cm}
\dom k = \left[0, +\infty\right[ & Zoals bij $\sqrt{x}$. \\
\ber k = \left]-\infty, 0\right] & Het bereik klapt om.
\end{stappenvrij}}{%
\begin{tikzpicture}[x=.45cm,y=.45cm,font=\scriptsize]
  \rooster{-1}{1}{4}{-2}{1}{2}
  \draw[grijs,dashed,thick] plot[domain=0:4,samples=60,smooth] (\x,{sqrt(\x)});
  \kromme{0}{4}{-sqrt(\x)}
  \vol{(0,0)}
\end{tikzpicture}}}

\subsection{Spiegelen om de $y$-as}

Met $b = -1$ verandert binnen de haakjes $x$ in $-x$: wat $f$ in $2$ doet, doet de nieuwe functie in $-2$.

\begin{theorie}
\begin{stelling}{Spiegelen om de $y$-as}
De grafiek van $\formule{h(x) = f(-x)}$ is het spiegelbeeld van die van $f$ om de $y$-as: \quad $\formule{(x, y) \mapsto (-x, y)}$
\end{stelling}

\begin{bewijs}
Het bewijs van horizontaal uitrekken gebruikt alleen dat $b \neq 0$. Met $b = -1$ geeft het: bij elk punt $(u, y)$ van $f$ hoort het punt $(-u, y)$ van $h$. Dat is het spiegelbeeld om de $y$-as.
\end{bewijs}
\end{theorie}

Nu klapt het domein om, en blijft het bereik. Dit sluit aan bij hoofdstuk 1: voor een even functie is $f(-x) = f(x)$, en spiegelen om de $y$-as verandert de grafiek niet. Voor een oneven functie is $f(-x) = -f(x)$: beide spiegelingen geven dan hetzelfde.

\voorbeeld{}{Geef het domein en het bereik van $h(x) = \sqrt{-x}$.\label{vb:h2-spiegely}}{%
\naast{%
\uitleg{Spiegel $\sqrt{x}$ om de $y$-as: de $y$-coördinaten blijven, de $x$-coördinaten veranderen van teken.}
\begin{stappenvrij}{4.2cm}{4.2cm}
\dom h = \left]-\infty, 0\right] & Het domein klapt om. \\
\ber h = \left[0, +\infty\right[ & Zoals bij $\sqrt{x}$.
\end{stappenvrij}
Controle: $h(-4) = \sqrt{4} = 2$, en $(-4, 2)$ is het spiegelbeeld van $(4, 2)$.}{%
\begin{tikzpicture}[x=.45cm,y=.45cm,font=\scriptsize]
  \rooster{-4}{1}{4}{-1}{1}{2}
  \draw[grijs,dashed,thick] plot[domain=0:4,samples=60,smooth] (\x,{sqrt(\x)});
  \kromme{-4}{0}{sqrt(-\x)}
  \vol{(0,0)}
\end{tikzpicture}}}

\nujij{Welke functie krijg je als je de grafiek van $x^3$ om de $y$-as spiegelt? En om de $x$-as?}{Om de $y$-as: $(-x)^3 = -x^3$. Om de $x$-as: $-x^3$. Hetzelfde, want $x^3$ is oneven.}

\begin{oefeningen}
\oef{h2-rekken}

\oef{h2-spiegelen}
\end{oefeningen}

%========================================================================================
\section{Transformaties combineren}

Bij een voorschrift als $g(x) = a \cdot f\bigl(b(x - c)\bigr) + d$ voer je de transformaties na elkaar uit. Welke volgorde is de goede? Volg een punt van $f$ en kijk waar het terechtkomt.

\begin{theorie}
\begin{stelling}{Een punt volgen}
Is $\formule{g(x) = a \cdot f\bigl(b(x - c)\bigr) + d}$ met $a, b \neq 0$, dan hoort bij elk punt $(u, v)$ van de grafiek van $f$ het punt \ $\formule{\bigl(\tfrac{u}{b} + c,\ av + d\bigr)}$ \ van de grafiek van $g$, en omgekeerd.
\end{stelling}

\begin{bewijs}
Neem een punt $(x, y)$ en noem $u = b(x - c)$ en $v = \frac{y - d}{a}$, zodat $x = \frac{u}{b} + c$ en $y = av + d$. Dan is
\bewijsregel{y = a f\bigl(b(x - c)\bigr) + d \iff av + d = a f(u) + d \iff v = f(u) .}
Het punt $(x, y)$ ligt dus op de grafiek van $g$ als en slechts als $(u, v)$ op die van $f$ ligt.
\end{bewijs}
\end{theorie}

Lees de formule als een recept. Voor de $x$-coördinaat deel je eerst door $b$ en tel je daarna $c$ op: eerst horizontaal uitrekken, dan verschuiven. Voor de $y$-coördinaat vermenigvuldig je eerst met $a$ en tel je daarna $d$ op: eerst verticaal uitrekken, dan verschuiven. Werk dus van binnen naar buiten.

\voorbeeld{}{Schets $g(x) = -2\sqrt{x + 1} + 3$ vanuit $\sqrt{x}$, en geef het domein en het bereik.\label{vb:h2-combineren}}{%
\naast{%
\uitleg{Volg het beginpunt $(0, 0)$ door alle stappen.}
\begin{stappenvrij}{3.9cm}{3.7cm}
\sqrt{x + 1} & 1 naar links: $(-1, 0)$. \\
2\sqrt{x + 1} & Verticaal $\times 2$: $(-1, 0)$. \\
-2\sqrt{x + 1} & Spiegelen om de $x$-as. \\
-2\sqrt{x + 1} + 3 & 3 naar boven: $(-1, 3)$. \\
\dom g = \left[-1, +\infty\right[ & \\
\ber g = \left]-\infty, 3\right] & De grafiek daalt vanaf $3$.
\end{stappenvrij}}{%
\begin{tikzpicture}[x=.5cm,y=.5cm,font=\scriptsize]
  \rooster{-1}{1}{5}{-2}{1}{3}
  \draw[grijs,dashed,thick] plot[domain=0:5,samples=60,smooth] (\x,{sqrt(\x)});
  \kromme{-1}{5}{-2*sqrt(\x+1)+3}
  \vol{(-1,3)}
\end{tikzpicture}}}

\begin{werkwijze}
Schrijf het voorschrift als $a \cdot f\bigl(b(x - c)\bigr) + d$. Pas toe in deze volgorde: horizontaal uitrekken met $\frac1b$ (en spiegelen als $b < 0$), verschuiven over $c$ naar rechts, verticaal uitrekken met $a$ (en spiegelen als $a < 0$), verschuiven over $d$ naar boven. Volg een bijzonder punt, en controleer met één functiewaarde.
\end{werkwijze}

\begin{oefeningen}
\oef{h2-combineren-schets}
\end{oefeningen}

Staat er een factor bij $x$, zet die dan eerst buiten haakjes. Anders lees je de verschuiving verkeerd af.

\voorbeeld{}{Schets $g(x) = \sqrt{2x - 6}$ vanuit $\sqrt{x}$.\label{vb:h2-buitenhaakjes}}{%
\naast{%
\uitleg{Een vlugge lezing zegt: 6 naar rechts. Maar dan zou het beginpunt $(6, 0)$ zijn, terwijl $g(6) = \sqrt6 \neq 0$. Zet eerst $2$ buiten haakjes.}
\begin{stappenvrij}{4.2cm}{4.4cm}
g(x) = \sqrt{2(x - 3)} & $b = 2$ en $c = 3$. \\
(0, 0) \mapsto (3, 0) & $\frac02 + 3 = 3$. \\
(4, 2) \mapsto (5, 2) & $\frac42 + 3 = 5$. \\
g(5) = \sqrt{4} = 2 & Controle.
\end{stappenvrij}
Eerst horizontaal $\times \frac12$, dan $3$ (niet $6$) naar rechts.}{%
\begin{tikzpicture}[x=.45cm,y=.45cm,font=\scriptsize]
  \rooster{-1}{1}{6}{-1}{1}{3}
  \draw[grijs,dashed,thick] plot[domain=0:6,samples=60,smooth] (\x,{sqrt(\x)});
  \kromme{3}{6}{sqrt(2*\x-6)}
  \vol{(3,0)} \vol{(5,2)}
\end{tikzpicture}}}

\begin{oefeningen}
\oef{h2-combineren}
\end{oefeningen}

Omgekeerd kun je uit een grafiek het voorschrift halen, als je weet van welke basisfunctie ze komt. Een bijzonder punt geeft $c$ en $d$, een tweede punt de factor $a$.

\voorbeeld{}{De parabool hiernaast heeft top $(-1, 2)$ en gaat door $(1, 0)$. Geef het voorschrift.\label{vb:h2-voorschrift}}{%
\naast{%
\uitleg{Een parabool is $x^2$, verschoven naar de top en verticaal uitgerekt: $y = a(x - c)^2 + d$. Bepaal $a$ met het tweede punt.}
\begin{stappenvrij}{4.6cm}{3.6cm}
y = a(x + 1)^2 + 2 & Top $(-1, 2)$. \\
0 = a \cdot 2^2 + 2 \Rightarrow a = -\tfrac12 & Punt $(1, 0)$ invullen. \\
y = -\tfrac12(x + 1)^2 + 2 &
\end{stappenvrij}}{%
\begin{tikzpicture}[x=.5cm,y=.5cm,font=\scriptsize]
  \rooster{-4}{1}{2}{-1}{1}{3}
  \kromme{-3.5}{1.5}{-(\x+1)^2/2+2}
  \vol{(-1,2)} \vol{(1,0)}
\end{tikzpicture}}}

\begin{oefeningen}
\oef{h2-voorschrift-kort}
\end{oefeningen}

%========================================================================================
\opdrachtenkop

\begin{oefeningen}
\needspace{10\baselineskip}
\oef{h2-o1}

\needspace{9\baselineskip}
\oef{h2-o10}

\oef{h2-o2}

\needspace{9\baselineskip}
\oef{h2-o11}

\needspace{7\baselineskip}
\oef{h2-o14}

\needspace{6\baselineskip}
\oef{h2-o12}

\needspace{8\baselineskip}
\oef{h2-o13}

\needspace{10\baselineskip}
\oef{h2-o16}

\needspace{9\baselineskip}
\oef{h2-o3}

\needspace{7\baselineskip}
\oef{h2-o20}

\needspace{8\baselineskip}
\oef{h2-o4}

\needspace{10\baselineskip}
\oef{h2-o6}

\needspace{9\baselineskip}
\oef{h2-o17}

\needspace{11\baselineskip}
\oef{h2-o5}

\needspace{9\baselineskip}
\oef{h2-o21}

\needspace{8\baselineskip}
\oef{h2-o24}

\needspace{9\baselineskip}
\oef{h2-o18}

\needspace{10\baselineskip}
\oef{h2-o25}

\needspace{10\baselineskip}
\oef{h2-o7}

\needspace{10\baselineskip}
\oef{h2-o27}

\needspace{8\baselineskip}
\oef{h2-o26}

\oef{h2-o8}

\needspace{8\baselineskip}
\oef{h2-o28}

\oef{h2-o9}
\end{oefeningen}

%----------------------------------------------------------------------------------------
% KOPPEL: verbanden met andere leerstof
%----------------------------------------------------------------------------------------
\needspace{20\baselineskip}
\opdrachtsectie{Koppel}

\lkr{Goniometrische functies, Exponentiële en logaritmische functies en Statistiek: kijk na wat al gezien is. Het pakket legt de volgorde niet vast.}

\begin{oefeningen}
\oef{h2-koppel-1}

\oef{h2-koppel-2}

\oef{h2-koppel-3}
\end{oefeningen}

%----------------------------------------------------------------------------------------
\opdrachtsectie{Test jezelf}

\begin{oefeningen}
\oef{h2-t5}

\oef{h2-t1}

\oef{h2-t2}

\oef{h2-t6}

\oef{h2-t3}

\oef{h2-t7}

\oef{h2-t4}
\end{oefeningen}

\herhaalschema{%
  \fzv{oef:h2-t5}{oef:h2-basis}%
  \fv{oef:h2-t1}{vb:h2-verschuiven}{oef:h2-verschuiven}%
  \fv{oef:h2-t2}{vb:h2-rekken}{oef:h2-rekken}%
  \fv{oef:h2-t6}{vb:h2-rekken}{oef:h2-o10}%
  \fv{oef:h2-t3}{vb:h2-combineren}{oef:h2-combineren}%
  \fv{oef:h2-t7}{vb:h2-combineren}{oef:h2-combineren-schets}%
  \fv{oef:h2-t4}{vb:h2-voorschrift}{oef:h2-voorschrift-kort}}

\ifzelfstudie
\section*{Maak je eigen samenvatting}
Sluit dit boek en schrijf in je schrift: de grafieken van de vijf basisfuncties met domein en bereik, wat $f(x - c) + d$, $a f(x)$, $f(bx)$, $-f(x)$ en $f(-x)$ met een grafiek doen, en in welke volgorde je combineert. Vergelijk daarna met de kaders.
\fi
